%% ========================================================================
%% Lampiran D: Glosarium
%% ========================================================================
%% STATUS: SCAFFOLD ONLY - definisi di bawah adalah stub satu baris; lihat
%% docs/authoring-guide.md
%% ========================================================================

\chapter{Glosarium}
\label{app:glosarium}

Daftar istilah kunci yang dipakai sepanjang buku ini, disusun menurut abjad.

\begin{description}[leftmargin=0pt, style=nextline]
\gitem{Artisan}{TODO: definisi ringkas CLI bawaan Laravel.}

\gitem{Blade}{TODO: definisi ringkas mesin templating Laravel.}

\gitem{Controller}{TODO: definisi ringkas kelas penangan request pada pola MVC.}

\gitem{Eager loading}{TODO: definisi ringkas teknik memuat relasi di muka untuk menghindari N+1 query.}

\gitem{Eloquent}{TODO: definisi ringkas ORM bawaan Laravel.}

\gitem{FormRequest}{TODO: definisi ringkas kelas validasi input terpisah dari controller.}

\gitem{Middleware}{TODO: definisi ringkas lapisan penyaring request/response.}

\gitem{Migrasi}{TODO: definisi ringkas berkas versi terkontrol untuk skema basis data.}

\gitem{Model}{TODO: definisi ringkas representasi tabel basis data sebagai kelas PHP.}

\gitem{RBAC}{TODO: definisi ringkas Role-Based Access Control.}

\gitem{Route}{TODO: definisi ringkas pemetaan URL ke controller/handler.}

\gitem{Sanctum}{TODO: definisi ringkas paket autentikasi token Laravel.}

\gitem{Seeder}{TODO: definisi ringkas mekanisme pengisi data awal/dummy basis data.}

\gitem{Migrasi rollback}{TODO: definisi ringkas pembatalan batch migrasi terakhir.}

\gitem{Token}{TODO: definisi ringkas kredensial akses API berumur terbatas.}
\end{description}
